\newacronym{gnb}{GNB}{Grupo de Neurocomputación Biológica}
\newacronym{uam}{UAM}{Universidad Autónoma de Madrid}


\section{Modelos de dinámica neuronal usados\label{SEC:MODELOS_DINAMICA_NEURONAL}}

En esta sección se abordan las dos piezas fundamentales del circuito biohíbrido: el modelo neuronal y el modelo sináptico. Es importante establecer una distinción entre ambos. Por un lado, el modelo neuronal de Hindmarsh-Rose ha sido escogido entre diversas alternativas, decisión que se justificará a continuación. Por otro lado, el modelo de sinapsis química no representa una simple elección de diseño, sino que el estudio y parametrización del mismo constituye el objeto central de este trabajo. En los siguientes apartados se describe en profundidad la formulación de ambos sistemas.

\subsection{Modelo de Hindmarsh-Rose modificado\label{SUBSEC:HINDMARSH_ROSE}}

Para el desarrollo de este trabajo se emplea una versión modificada del modelo de Hindmarsh-Rose \cite{Hindmarsh1984}. Aunque constituye una descripción matemática simplificada en comparación con otras formulaciones biofísicas de mayor detalle, el modelo posee la fidelidad suficiente para los requerimientos experimentales de este proyecto, reproduciendo con precisión patrones complejos como el disparo en ráfagas (\textit{bursting}).

El sistema dinámico continuo se define mediante el conjunto de ecuaciones diferenciales ordinarias de la Ecuación~\ref{EQ:HR_SYSTEM}:

\begin{equation}[EQ:HR_SYSTEM]{Sistema de ecuaciones del modelo de Hindmarsh-Rose modificado}
    \boxed{
        \begin{aligned}
            \frac{dx(t)}{dt}               & = y(t) + b x^2(t) - a x^3(t) - z(t) + e - I_{\mathrm{syn}}(t) \\
            \frac{dy(t)}{dt}               & = c - d x^2(t) - y(t)                                         \\
            \frac{1}{\mu} \frac{dz(t)}{dt} & = -v_h \cdot z(t) + S(x(t) - x_r)
        \end{aligned}
    }
\end{equation}

En este sistema, la variable \(x(t)\) representa el potencial de membrana adimensional. En la parametrización empleada, esta variable se desplaza típicamente en una escala comprendida entre \(-1{,}6\) y \(1{,}7\), por lo que tanto dichos valores como la propia variable se expresan en unidades adimensionales. Además, el término \(I_{\mathrm{syn}}(t)\) representa la corriente sináptica aplicada a la neurona, que actúa como una entrada externa adicional sobre la dinámica rápida del potencial de membrana (no es un parámetro). Es importante destacar que esta corriente se incluye con signo negativo en el sistema; esta convención garantiza que, bajo la formulación del modelo de sinapsis descrito en la Sección~\ref{SUBSEC:SINAPSIS_QUIMICA}, un potencial de reversión \(E_{\mathrm{syn}}\) positivo (superior al potencial de membrana) ejerza un efecto excitatorio. Esto permite alinearse con la interpretación biofísica estándar y con la disposición de la corriente sináptica en el modelo neuronal (de diferente naturaleza) presentado en \cite{Golowasch1999}. Por su parte, \(y(t)\) actúa como la variable de disparo rápido que cuantifica la tasa de transporte de iones de sodio y potasio a través de los canales iónicos rápidos de membrana, gobernando la dinámica de recuperación del potencial de acción, y \(z(t)\) constituye la corriente de adaptación lenta (asociada específicamente al transporte de iones lentos), encargada de incrementarse con cada potencial de acción y reducir progresivamente la tasa de disparo, modulando así la alternancia entre episodios de actividad y períodos de silencio.

Con el objetivo de asegurar un comportamiento rítmico consistente en el circuito biohíbrido, se han adoptado dos modificaciones paramétricas empleadas por el \ac{gnb} de la \ac{uam}. Por un lado, al establecer $v_h = 0{,}1$ en lugar del valor unitario presente en el modelo original, se reduce el decaimiento natural de la variable de adaptación lenta. Por otro lado, la reducción de la constante de acoplamiento de su valor canónico $S = 4{,}0$ a $S = 1{,}0$ atenúa la influencia del voltaje sobre dicha variable lenta. Esta configuración asegura de forma conjunta una elevada estabilidad rítmica en las neuronas, limitando la aparición de dinámicas caóticas y estabilizando el patrón de ráfagas para facilitar una interacción biohíbrida robusta.

El resto de los parámetros se han seleccionado conforme a los valores habitualmente utilizados en la literatura matemática para el modelo. Cada uno de estos coeficientes tiene una función específica en la geometría del atractor dinámico, garantizando el comportamiento estable requerido. El comportamiento resultante de la neurona utilizando este conjunto completo de parámetros se ilustra en la Figura~\ref{FIG:HR_ROSE}. El significado fenomenológico y los valores exactos empleados del resto de parámetros se detallan a continuación:

\begin{itemize}
    \item $e = 3{,}281$: Corriente de estímulo externo que modula el régimen de disparo y facilita la aparición de ráfagas.
    \item $\mu = 0{,}0021$: Parámetro de escala temporal que gobierna la lentitud intrínseca de la variable de adaptación $z(t)$.
    \item $a = 1{,}0$: Coeficiente cúbico de la dinámica rápida que modula la amplitud y forma de los potenciales de acción.
    \item $b = 3{,}0$: Coeficiente cuadrático de la dinámica rápida que regula la excitabilidad de membrana e influye en la transición entre ráfagas y disparo tónico.
    \item $c = 1{,}0$: Término constante que establece el valor de la corriente de recuperación rápida $y(t)$ en ausencia de voltaje.
    \item $d = 5{,}0$: Coeficiente cuadrático que determina cómo el voltaje $x(t)$ activa la corriente de recuperación $y(t)$.
    \item $x_r = -1{,}6$: Potencial de referencia que define el estado de reposo relativo del subsistema lento.
\end{itemize}

\begin{figure}[Dinámica de la neurona de Hindmarsh-Rose modificada]{FIG:HR_ROSE}{Simulación temporal del modelo de Hindmarsh-Rose modificado. El ajuste conjunto de $v_h$ y $S$ junto con el resto de la parametrización estándar permite obtener un patrón de actividad con una elevada estabilidad rítmica.}
    \image{0.9\textwidth}{}{hindmarsh-rose_modified}
\end{figure}


\subsection{Modelo de sinapsis química\label{SUBSEC:SINAPSIS_QUIMICA}}

El modelo de sinapsis química empleado se fundamenta en dos componentes que actúan en paralelo: una corriente rápida (\(I_{\mathrm{fast}}\)), que en este trabajo se ha vinculado únicamente a las altas frecuencias o \textit{spikes} de \(V_{\mathrm{pre}}\) pese a poder modelar ambas, como se justifica en la Sección~\ref{SUBSEC:USO_SIMULTANEO}; y una corriente lenta (\(I_{\mathrm{slow}}\)), vinculada a las bajas frecuencias o \textit{bursts}. Ambas contribuciones dependen de los potenciales presináptico (\(V_{\mathrm{pre}}\)) y postsináptico (\(V_{\mathrm{post}}\)) conforme a la formulación matemática de \cite{Golowasch1999}, de donde también se adoptan las unidades físicas usadas en esta sección.

\subsubsection{Componente rápida (\(I_{\mathrm{fast}}\))}

La componente rápida es una corriente graduada instantánea que depende de \(V_{\mathrm{pre}}(t)\) mediante una sigmoide (cuya salida está acotada entre 0 y 1). Se define mediante la Ecuación~\ref{EQ:IFAST}:
\begin{equation}[EQ:IFAST]{Fórmula de la corriente rápida de la sinapsis química}
    \boxed{
        I_{\mathrm{fast}}(t) =
        g_{\mathrm{fast}}
        \cdot
        \frac{V_{\mathrm{post}}(t) - E_{\mathrm{syn}}}
        {1 + \exp[ s_{\mathrm{fast}} \cdot (V_{\mathrm{fast}} - V_{\mathrm{pre}}(t)) ]}
    }
\end{equation}
Los parámetros que aparecen en esta expresión son:
\begin{itemize}
    \item \(g_{\mathrm{fast}}\) (\(\mathrm{\mu S}\)): conductancia máxima rápida. Representa la densidad efectiva de canales sinápticos rápidos disponibles y actúa como un factor de ganancia que modifica la escala (amplitud) de la curva de corriente generada, pero sin alterar su morfología temporal; valores mayores permiten que \(I_{\mathrm{fast}}\) tenga un impacto más marcado sobre \(V_{\mathrm{post}}\) (mayor fuerza de la sinapsis), debiendo ajustarse según lo que se busque en este aspecto.
    \item \(E_{\mathrm{syn}}\) (mV): potencial de reversión sináptico. El hecho de que la sinapsis actúe de forma excitatoria o inhibitoria radica en este parámetro, ya que al evaluar el término \((V_{\mathrm{post}} - E_{\mathrm{syn}})\) determina el signo final de la corriente. Dependiendo de si se vuelve negativa o positiva, la corriente inyectada a la neurona postsináptica tendrá un efecto distinto: si \(E_{\mathrm{syn}}\) se sitúa claramente por encima del rango típico de \(V_{\mathrm{post}}\), la corriente inyectada neta es negativa, por lo que la influencia de \(V_{\mathrm{pre}}\) es directa (despolarizante al restarse en las ecuaciones del modelo), comportándose de forma excitatoria y tendiendo a favorecer patrones en fase. Si se coloca claramente por debajo de dicho rango, la corriente inyectada se vuelve positiva, lo que invierte la influencia de \(V_{\mathrm{pre}}\) (reflexión sobre el eje horizontal), actuando de forma inhibitoria (hiperpolarizante) y favoreciendo patrones en antifase. Además, cuanto más alejado se encuentre \(E_{\mathrm{syn}}\) del rango de operación de \(V_{\mathrm{post}}\), menor será la influencia dinámica de \(V_{\mathrm{post}}\) frente a \(V_{\mathrm{pre}}\) sobre la forma de la curva resultante. El hecho de que la morfología de la corriente replique o invierta la de \(V_{\mathrm{pre}}\) facilite la consecución de patrones en fase o antifase se debe a que garantiza que la sinapsis excite o inhiba a la neurona postsináptica precisamente cuando la presináptica dispara y viceversa cuando se hiperpolariza. Esta mitigación de la influencia postsináptica es, por tanto, de gran interés para forzar un acoplamiento dominante, siempre manteniendo la cautela de no exceder los rangos fisiológicamente plausibles para el tipo de sinapsis modelada. Además, este parámetro afectará también, como es lógico, a la escala de la corriente, aumentándola cuanto más se aleje de los límites de \(V_{\mathrm{post}}\), aunque esto puede ser compensado con \(g_{\mathrm{fast}}\).
    \item \(V_{\mathrm{pre}}(t)\), \(V_{\mathrm{post}}(t)\) (mV): potenciales de membrana pre y postsináptico.
    \item \(s_{\mathrm{fast}}\) (\(\mathrm{mV}^{-1}\)): pendiente de la sigmoide rápida. Valores grandes producen una transición muy abrupta entre estado inactivo y activo; esto implica que, cuanto mayor sea su valor, más compactación vertical sufrirá la parte de \(V_{\mathrm{pre}}\) que caiga en los bordes de la sigmoide, y más se expandirá verticalmente la parte que caiga en el centro. Por el contrario, valores pequeños generan una dependencia más suave de la corriente respecto a \(V_{\mathrm{pre}}\).
    \item \(V_{\mathrm{fast}}\) (mV): voltaje de semiactivación rápida (umbral de la sigmoide). Constituye el punto matemático en el que la función sigmoide evalúa exactamente a 0.5, por lo que actúa como punto de centraje matemático sobre \(V_{\mathrm{pre}}\).
\end{itemize}

Ambos parámetros (\(s_{\mathrm{fast}}\) y \(V_{\mathrm{fast}}\)) se combinan para lograr el comportamiento deseado de la rama rápida. Al posicionar \(V_{\mathrm{fast}}\) en el centro de los \textit{spikes} y aplicar un \(s_{\mathrm{fast}}\) grande, se consigue que los picos ocupen todo el ancho activo de la sigmoide. Simultáneamente, esta combinación aplanará la onda lenta (los \textit{bursts}) en el extremo inferior de la curva, suprimiéndola casi por completo. Como resultado, \(I_{\mathrm{fast}}\) genera pulsos estrechos fielmente sincronizados con los \textit{spikes} o componentes de alta frecuencia, evitando reflejar la componente lenta de la señal presináptica.

\subsubsection{Componente lenta (\(I_{\mathrm{slow}}\))}

La componente lenta introduce una compuerta dinámica \(m_{\mathrm{slow}}(t)\) que integra la historia reciente de \(V_{\mathrm{pre}}\). El sistema que define tanto la corriente como la evolución de dicha compuerta se detalla en la Ecuación~\ref{EQ:ISLOW_SYSTEM}:

\begin{equation}[EQ:ISLOW_SYSTEM]{Sistema de ecuaciones de la corriente lenta de la sinapsis química}
    \boxed{
        \begin{aligned}
            I_{\mathrm{slow}}(t)              & = g_{\mathrm{slow}} \cdot m_{\mathrm{slow}}(t) \cdot (V_{\mathrm{post}}(t) - E_{\mathrm{syn}})                                                                \\
            \frac{d m_{\mathrm{slow}}(t)}{dt} & = \frac{k_1 \cdot (1 - m_{\mathrm{slow}}(t))}{1 + \exp[ s_{\mathrm{slow}} \cdot (V_{\mathrm{slow}} - V_{\mathrm{pre}}(t)) ]} - k_2 \cdot m_{\mathrm{slow}}(t)
        \end{aligned}
    }
\end{equation}

donde \(g_{\mathrm{slow}}\) (\(\mathrm{\mu S}\)) y \(E_{\mathrm{syn}}\) (mV) funcionan igual que en la versión rápida, pero para la componente lenta, y \(m_{\mathrm{slow}}(t)\) es una variable adimensional que representa la fracción efectiva de canales lentos abiertos. Se puede apreciar, que ahora \(V_{\mathrm{pre}}\) afecta a la corriente a través de la sigmoide y esta, a su vez, mediante la variable \(m_{\mathrm{slow}}\); en lugar de únicamente con la sigmoide como ocurría en la componente rápida.

La dinámica de \(m_{\mathrm{slow}}(t)\) garantiza matemáticamente su acotación entre 0 y 1 gracias al compromiso dinámico entre la apertura impulsada por \(k_1\) (cuyo término se anula si \(m_{\mathrm{slow}} = 1\)) y el cierre impulsado por \(k_2\) (que se anula si \(m_{\mathrm{slow}} = 0\)), con:
\begin{itemize}
    \item \(k_1\) (\(\mathrm{ms}^{-1}\)): tasa máxima de apertura lenta. Su magnitud determina la inercia temporal de la compuerta y depende directamente de la frecuencia de corte que separa la dinámica rápida de los \textit{spikes} y la dinámica lenta del \textit{burst}. Valores mayores (constantes de tiempo menores) hacen que \(m_{\mathrm{slow}}\) siga con poca inercia las despolarizaciones permitiendo capturar respuestas de mayor agilidad en la compuerta (siempre que exista una separación temporal con la escala de los \textit{spikes}), mientras que valores menores (constantes comparables o mayores que el periodo de los \textit{bursts}) enfatizan la integración de la dinámica lenta y lo convierten en un integrador lento.
    \item \(k_2\) (\(\mathrm{ms}^{-1}\)): tasa de cierre, expresada también empíricamente. Controla la duración de la compuerta abierta tras cesar la despolarización; valores relativamente pequeños respecto a \(k_1\) permiten que \(m_{\mathrm{slow}}\) permanezca elevado a lo largo del \textit{burst} y decaiga entre \textit{bursts}.
    \item \(s_{\mathrm{slow}}\) (\(\mathrm{mV}^{-1}\)) y \(V_{\mathrm{slow}}\) (mV): pendiente y voltaje de semiactivación lenta. Funcionan igual que en la versión rápida, pero para la componente lenta, afectando ahora a través de \(m_{\mathrm{slow}}\), modulada por \(k_1\) y \(k_2\), en lugar de directamente.
\end{itemize}

De manera análoga, \(s_{\mathrm{slow}}\) y \(V_{\mathrm{slow}}\) se combinan para aislar la envolvente del \textit{burst}. Dado que los \textit{spikes} de alta frecuencia ya quedan cortados y mitigados por la inercia temporal impuesta por \(k_1\) y \(k_2\), se busca un valor de \(s_{\mathrm{slow}}\) pequeño (sin pasarse, para evitar que el rango dinámico de \(m_{\mathrm{slow}}\) sea minúsculo) que prevenga la deformación de la onda lenta. Al centrar \(V_{\mathrm{slow}}\) respecto a esta onda lenta y emplear dicha pendiente reducida, se logra que el \textit{burst} ocupe todo el ancho de la sigmoide sin distorsionarse. Con esta configuración, \(I_{\mathrm{slow}}\) actúa como una representación fidedigna en baja frecuencia de la señal presináptica.

\subsubsection{Uso simultáneo de corrientes\label{SUBSEC:USO_SIMULTANEO}}

En ocasiones resulta conveniente estudiar dinámicas aisladas empleando únicamente la corriente rápida \(I_{\mathrm{fast}}\), centrada en los \textit{spikes} de \(V_{\mathrm{pre}}\); o la lenta \(I_{\mathrm{slow}}\), centrada en los \textit{bursts} de \(V_{\mathrm{pre}}\). No obstante, si bien en la referencia original \cite{Golowasch1999} estas componentes se tratan de forma independiente, la arquitectura del repositorio del Apéndice~\ref{SEC:NEUNREPO} permite integrar la contribución de ambas en una única corriente sináptica total (Ecuación~\ref{EQ:ISYN}), logrando una reconstrucción sináptica completa de la forma de la señal presináptica:

\begin{equation}[EQ:ISYN]{Fórmula de la corriente total de la sinapsis química}
    \boxed{ I_{\mathrm{syn}}(t) = I_{\mathrm{fast}}(t) + I_{\mathrm{slow}}(t) }
\end{equation}

donde \(I_{\mathrm{syn}}\) se expresa en unidades de corriente (nA). Es importante notar que \(E_{\mathrm{syn}}\) será el único parámetro compartido por ambas componentes.

Algo se podría pensar es que \(I_{\mathrm{fast}}\), con un \(s_{\mathrm{fast}}\) pequeño para que no se deforme la influencia de \(V_{\mathrm{pre}}(t)\), y \(V_{\mathrm{fast}}\) adecuadamente centrado al centro del rango de \(V_{\mathrm{pre}}\), captura intrínsecamente todas las frecuencias presinápticas, por lo que teóricamente bastaría con esta única corriente para replicar la señal presináptica al completo sin requerir de esta reconstrucción sumando. Esto es completamente cierto, pero en este diseño se ha optado deliberadamente por restringir \(I_{\mathrm{fast}}\) a las altas frecuencias con un propósito fundamental: la modularidad. Disgregar la señal en dos vías frecuenciales permite escalar sus componentes de forma independiente. Partiendo de un estado base que replique \(V_{\mathrm{pre}}\) mediante ambas componentes sumadas (equilibrando \(g_{\mathrm{fast}}\) y \(g_{\mathrm{slow}}\)), se tendrá la flexibilidad de, por ejemplo, amplificar o atenuar los \textit{spikes} sin alterar los \textit{bursts} y viceversa. Esta capacidad de modulación es interesante para diversos experimentos, entre ellos los desarrollados por el \ac{gnb} de la \ac{uam}.

\subsubsection{Interacciones bidireccionales}

En interacciones bidireccionales se consideran dos sinapsis del tipo descrito, una en cada dirección: cada una utiliza como \(V_{\mathrm{pre}}\) el potencial de una de las neuronas e inyecta su corriente \(I_{\mathrm{syn}}\) en la otra, con sus propios parámetros y con \(V_{\mathrm{post}}\) definido en el compartimento somático correspondiente.


\section{Optimización automática de los parámetros de la sinapsis química\label{SEC:OPTIMIZACION_AUTOMATICA}}

Para que el modelo sináptico descrito en la Sección~\ref{SUBSEC:SINAPSIS_QUIMICA} reproduzca una dinámica objetivo concreta es necesario encontrar la combinación de parámetros que minimice la discrepancia entre la corriente sináptica generada y el comportamiento deseado. Esto requiere dos elementos fundamentales: una función de evaluación que cuantifique la calidad de cada combinación de parámetros, y un algoritmo de optimización que explore el espacio paramétrico de forma eficiente.

Un requisito de diseño central es que el algoritmo de parametrización no sea de un solo uso, sino que pueda aplicarse de forma recurrente en cada experimento con neuronas distintas, adaptándose a las características electrofisiológicas particulares de cada preparación biológica. En consecuencia, el método de optimización debe ser eficiente en el número de evaluaciones de la función objetivo, ya que cada evaluación resulta costosa en tiempo, especialmente en experimentos \textit{online} (con neuronas vivas y no con registros en frío) en los que no se dispone de tiempo ilimitado, que son los de mayor interés. Debido a esto y a otros factores, se ha optado por utilizar la optimización bayesiana.

Se han desarrollado dos implementaciones complementarias de esta estrategia. Por un lado, una versión \textit{offline} que parametriza la sinapsis a partir de un registro de voltaje presináptico almacenado en un fichero; esta modalidad resulta muy rápida al no depender de los tiempos de adquisición de una neurona viva (el modelo sináptico se puede simular a velocidad arbitraria), aunque solo admite sinapsis unidireccional al disponer únicamente de una señal \(V_{\mathrm{pre}}\) estática. Por otro lado, una versión \textit{online} implementada como módulo de RTXI, que ejecuta la optimización en tiempo real durante el experimento, permitiendo su uso tanto en configuraciones unidireccionales como bidireccionales con neuronas vivas y modelo.

En las secciones siguientes se describe la función de evaluación y el algoritmo de optimización elegido, y luego se detallan las características de cada implementación.


\subsection{Función de evaluación\label{SUBSEC:FUNCION_EVALUACION}}

La función de evaluación es necesaria para poder cuantificar la calidad de una combinación de parámetros sinápticos midiendo cuánto se asemeja la corriente sináptica generada al comportamiento de referencia esperado. El algoritmo de evaluación desarrollado se descompone en dos componentes independientes: una puntuación de forma (\textit{shape score}) y una puntuación de rango (\textit{range score}), que capturan respectivamente la forma (sin contar escala y desplazamiento) y la amplitud de la corriente. En caso de usar ambas componentes (\(I_{\mathrm{fast}}\) e \(I_{\mathrm{slow}}\)), cada una se evalúa por separado y se promedian con pesos configurables (por defecto, 0{,}5 para cada una).

\subsubsection{Separación frecuencial}

Antes de calcular las puntuaciones, es necesario separar la señal presináptica \(V_{\mathrm{pre}}(t)\) en sus componentes de alta y baja frecuencia para generar las señales de referencia contra las que comparar \(I_{\mathrm{fast}}\) e \(I_{\mathrm{slow}}\). Esta separación se realiza mediante un filtro Butterworth paso bajo de cuarto orden \cite{Butterworth1930}, cuya función de transferencia se implementa en secciones de segundo orden (SOS) mediante la función \texttt{to\_sos} para maximizar la estabilidad numérica \cite{Proakis2007}. Para evitar desfases temporales, el filtro se aplica de forma bidireccional (fase cero) utilizando la función \texttt{filtfilt} \cite{Gustafsson1996} (Código~\ref{COD:BUTTERWORTH}). Además, tal y como aplican las implementaciones de referencia en procesamiento digital de señales \cite{SciPyfiltfilt}, resulta indispensable mitigar los transitorios o efectos de borde generados durante el filtrado. Para ello, se lleva a cabo una extrapolación de puntos en los extremos de la señal (relleno), técnica que consiste en añadir un margen temporal calculando la reflexión, en este caso impar, de los datos respecto a los puntos inicial y final, lo que garantiza la continuidad de la derivada en los extremos, resultando conveniente para evitar transitorios indeseados al aplicar el filtro (Código~\ref{COD:PADDING}). La frecuencia de corte \(f_c\) la define el usuario en función de la separación temporal deseada entre \textit{spikes} y \textit{bursts} (típicamente 0{,}3\,kHz para neuronas vivas, aunque puede variar entre las mismas o con las neuronas modelo). La componente filtrada constituye la referencia de baja frecuencia para \(I_{\mathrm{slow}}\), mientras que el residuo (señal original menos filtrada) constituye la referencia de alta frecuencia para \(I_{\mathrm{fast}}\).

Toda la lógica de procesamiento de señal se implementa mediante la biblioteca KFR (Apéndice~\ref{SUBSEC:KFRREPO}), seleccionada por tres motivos: (1)~la simplicidad de sus operaciones vectoriales sobre señales, que permite expresar operaciones como el centrado, la correlación o la extracción de extremos en pocas líneas de código; (2)~la eficiencia de dichas operaciones, optimizadas internamente mediante operaciones vectorizadas, vital en este caso dada la longitud en puntos de las señales (debido a las altas tasas de muestreo de estos experimentos); y (3)~la disponibilidad directa de funciones de diseño de filtros IIR (\texttt{iir\_lowpass}, \texttt{butterworth}, \texttt{to\_sos}) y de filtrado bidireccional (\texttt{filtfilt}), que simplifican notablemente la implementación del preprocesamiento frecuencial.

\subsubsection{Puntuación de forma}

La similitud en forma entre la corriente sináptica candidata y la referencia se evalúa mediante el coeficiente de correlación de Pearson \cite{Pearson1896} (Código~\ref{COD:PEARSON}). Dada una señal candidata \(\mathbf{x}\) y una señal de referencia \(\mathbf{y}\) de longitud \(n\), el coeficiente se calcula mediante la Ecuación~\ref{EQ:PEARSON}:
\begin{equation}[EQ:PEARSON]{Fórmula del coeficiente de correlación de Pearson}
    \boxed{
        r = \frac{\sum_{i=1}^{n} (x_i - \bar{x})(y_i - \bar{y})}{\sqrt{\sum_{i=1}^{n} (x_i - \bar{x})^2} \sqrt{\sum_{i=1}^{n} (y_i - \bar{y})^2}}
    }
\end{equation}
donde \(\bar{x}\) e \(\bar{y}\) representan las medias de las respectivas señales, y \(r \in [-1, 1]\). Este coeficiente se normaliza al intervalo \([0, 1]\) mediante la transformación \((r + 1) / 2\). De esta forma, una puntuación de 1 corresponde a una correlación perfecta (formas iguales), mientras que una puntuación de 0 corresponde a una correlación inversa (formas opuestas o simétricas respecto al eje horizontal), y todo esto sin tener en cuenta la escala y el desplazamiento de ninguna de las señales, únicamente su morfología (es adimensional).

Por tanto, la búsqueda de un acoplamiento en antifase se basa en una correlación directa, y la búsqueda en fase invierte esta puntuación para buscar una correlación inversa. Esto es de esta forma para compensar el hecho de que el modelo neuronal resta la corriente inyectada en sus ecuaciones.

\subsubsection{Puntuación de rango}

La puntuación de rango (Código~\ref{COD:RANGESCORE}) mide cuánto se desvían los extremos observados de la corriente candidata respecto a los buscados, complementando la puntuación de forma, que no captura información sobre el rango. Se calcula mediante la Ecuación~\ref{EQ:RANGESCORE} a continuación:
\begin{equation}[EQ:RANGESCORE]{Fórmula de la puntuación del rango de la corriente sináptica}
    \boxed{
    \mathrm{score}_{\mathrm{range}} = 1 - \frac{\frac{1}{2}\left(|I_{\min}^{\mathrm{obs}} - I_{\min}^{\mathrm{exp}}| + |I_{\max}^{\mathrm{obs}} - I_{\max}^{\mathrm{exp}}|\right)}{d_{\max}}
    }
\end{equation}
donde \(d_{\max}\) es una distancia máxima admisible que normaliza el error. Esta distancia se deriva del rango de la corriente esperada y se calcula mediante la Ecuación~\ref{EQ:DMAX}:
\begin{equation}[EQ:DMAX]{Fórmula de la distancia máxima para la normalización del error de la puntuación del rango}
    \boxed{
        \begin{aligned}
            R_{\mathrm{exp}} & = I_{\max}^{\mathrm{exp}} - I_{\min}^{\mathrm{exp}}                                                             \\
            d_{\max}         & = \left( R_{\mathrm{exp}} + R_{\mathrm{exp}} \cdot f_{\mathrm{margen}} \cdot 2 \right) \cdot f_{\mathrm{rango}}
        \end{aligned}
    }
\end{equation}
En esta expresión, \(d_{\max}\) representa la máxima diferencia de intensidad aproximada (porque realmente no está acotada) que habrá entre dos puntos. El parámetro \(f_{\mathrm{margen}}\) introduce unos márgenes de seguridad sobre el rango esperado \(R_{\mathrm{exp}}\), que es el mismo factor que se usará para definir posteriormente los límites de búsqueda de las conductancias para que no se llegue a corrientes extremas, lo que garantiza la coherencia del sistema. Por su parte, el factor empírico \(f_{\mathrm{rango}}\) se ha ajustado para asegurar que ambas componentes de la puntuación (rango y forma) presenten una variabilidad similar. De este modo, al no estar acotado estrictamente el error de rango (pudiendo resultar la puntuación final menor que cero en casos extremos), se evita que esta componente acabe dominando sobre la de forma y se preserva el peso de la morfología en la evaluación.

\paragraph{Asignación de rangos esperados: uso de una componente de la corriente frente a uso conjunto}

El tratamiento de los rangos esperados difiere según se utilice una única componente o ambas simultáneamente:

\begin{itemize}
    \item \textbf{Uso individual} (\(I_{\mathrm{fast}}\) o \(I_{\mathrm{slow}}\) por separado): el rango esperado de corriente es directamente el intervalo \([I_{\min}^{\mathrm{exp}}, I_{\max}^{\mathrm{exp}}]\) introducido por el usuario.
    \item \textbf{Uso conjunto} (\(I_{\mathrm{fast}}\) e \(I_{\mathrm{slow}}\) simultáneamente): el rango global introducido por el usuario se reparte proporcionalmente entre ambas componentes de forma coherente con un buen diseño de filtro paso bajo. En concreto, la componente lenta \(I_{\mathrm{slow}}\) hereda tanto el desplazamiento como la onda lenta de la señal presináptica (dado que el filtro paso bajo utilizado mantiene las magnitudes), por lo que su rango esperado se reescala con desplazamiento a partir del rango de \(V_{\mathrm{pre}}\) al rango de corriente total (Código~\ref{COD:RESCALE}, función \texttt{rescale\_to\_target}). La componente rápida \(I_{\mathrm{fast}}\), al ser el residuo tras la sustracción de la componente filtrada, está centrada en cero y carece de desplazamiento, por lo que su rango se reescala sin desplazamiento, aplicando únicamente la proporción entre rangos (función \texttt{rescale\_to\_target\_no\_offset}). Este reparto garantiza que la suma de ambas componentes reconstruya el rango total esperado, asignando a cada una la contribución proporcional a su peso frecuencial en la señal original.
\end{itemize}
